% =========================================================
\section{Objetivos y antecedentes}
% =========================================================

\begin{frame}{Objetivo general y objetivos específicos}
  \begin{UCBlock}{Objetivo general}
    Inserte aquí el objetivo general de la tesis de posgrado, formulado con un verbo claro y alineado con la pregunta de investigación.
  \end{UCBlock}
  \vspace{0.08cm}
  \begin{enumerate}
    \item Inserte aquí el primer objetivo específico.
    \item Inserte aquí el segundo objetivo específico.
    \item Inserte aquí el tercer objetivo específico.
    \item Inserte aquí el cuarto objetivo específico, si corresponde.
  \end{enumerate}
\end{frame}

\begin{frame}{Marco teórico o antecedentes}
  \begin{columns}[T,onlytextwidth]
    \begin{column}{0.32\textwidth}
      \begin{UCBlock}{Concepto 1}
        Inserte aquí una definición, teoría o referencia clave.
      \end{UCBlock}
    \end{column}
    \begin{column}{0.32\textwidth}
      \begin{UCBlock}{Concepto 2}
        Inserte aquí un antecedente metodológico o técnico relevante.
      \end{UCBlock}
    \end{column}
    \begin{column}{0.32\textwidth}
      \begin{UCBlock}{Concepto 3}
        Inserte aquí el antecedente que conecta con la brecha.
      \end{UCBlock}
    \end{column}
  \end{columns}
  \begin{UCKeyPoint}
    Evita convertir esta sección en una revisión bibliográfica extensa. Prioriza los conceptos que justifican el método y los objetivos.
  \end{UCKeyPoint}
\end{frame}

\begin{frame}{Alcance de la tesis de posgrado}
  \begin{columns}[T,onlytextwidth]
    \begin{column}{0.58\textwidth}
      \centering
      \begin{tikzpicture}[x=1cm,y=1cm]
        \node (a) at (0,3.12) {\UCProcessArrow{Entrada}{datos, muestras o caso}};
        \node (b) at (2.05,3.12) {\UCProcessArrow{Preparación}{limpieza o diseño}};
        \node (c) at (4.42,3.12) {\UCProcessArrow{Análisis}{modelo o ensayo}};
        \node (d) at (0.70,1.36) {\UCProcessArrow{Validación}{criterio definido}};
        \node (e) at (2.94,1.36) {\UCProcessArrow{Resultados}{indicadores clave}};
        \node (f) at (5.16,1.36) {\UCProcessArrow{Cierre}{discusión y límites}};
        \draw[-{Latex[length=2mm]},UCBluePPT,line width=0.7pt] (0.95,3.12)--(1.56,3.12);
        \draw[-{Latex[length=2mm]},UCBluePPT,line width=0.7pt] (3.02,3.12)--(3.78,3.12);
        \draw[-{Latex[length=2mm]},UCBluePPT,line width=0.7pt] (4.40,2.62)--(0.75,1.89);
        \draw[-{Latex[length=2mm]},UCBluePPT,line width=0.7pt] (4.40,2.62)--(2.92,1.89);
        \draw[-{Latex[length=2mm]},UCYellowDeep,line width=0.8pt] (3.62,1.36)--(4.22,1.36);
      \end{tikzpicture}
    \end{column}
    \begin{column}{0.38\textwidth}
      \begin{UCKeyPoint}
        Indica explícitamente qué queda dentro y fuera del alcance. Esto ayuda a controlar expectativas durante el avance.
      \end{UCKeyPoint}
    \end{column}
  \end{columns}
\end{frame}
